\documentclass[10pt]{article}
\usepackage{icmlko}

\icmlmeta{IP 파운데이션 월드모델}{975}{무엇이 capability 를 정하나 --- 요인 하나만 흔든 짝 실험}{어느 요인이 유보 예측 성능을 정하는지 묻는 자리가 그 동안 비어 있었다. 요인 다섯을 하나씩만 흔들어 델타와 짝지은 표준오차를 내고 순위를 적는다}

\begin{document}

\twocolumn[
\icmlheader

\begin{abstract}
「무엇이 capability 를 정하나」는 이 실험실의 여섯 능력 축 가운데 하나인데,
그 칸은 세 사이클 연속 \textbf{비어 있었다}. 자가 없어서가 아니다 ---
자는 이미 적혀 있었다(\textbf{요인 하나만 흔든 짝 실험의 $\Delta \pm SE$}).
없던 것은 \textbf{흔들 요인의 목록}이다. 흔들 요인을 등록한 적이 없으니
어떤 개선의 크기도 \textbf{견줄 대상이 없었다}.

이 노트는 그 목록을 먼저 만든다. 학습 파이프라인의 다섯 칸 ---
\textbf{architecture}(특징 집합) · \textbf{objective}(도메인 가중) ·
\textbf{data mixture}(원천 혼합) · \textbf{optimizer}(능형 정칙화) ·
\textbf{curriculum}(학습 창 길이) --- 에 격자를 \textbf{측정 전에} 못 박고,
밑판에서 \textbf{한 칸만} 바꾼 짝 실험을 돌린다.
유보 집합은 \textbf{모든 수준에서 글자 그대로 같다}.
$SE$ 는 유보 행을 도메인 안에서 되뽑는 \textbf{짝지은 군집 붓스트랩}에서 낸다.

결과는 순위 하나와 부호 하나다. 순위는 \textbf{특징 집합 $>$ 학습 창 $>$
원천 혼합 $>$ 도메인 가중 $>$ 능형 정칙화}이고, 🔴 \textbf{사전등록 문턱을 넘은
수준은 전부 $\Delta$ 가 음수다.} 즉 이 격자가 잰 것은
\textbf{「무엇이 성능을 올리나」가 아니라 「무엇을 망가뜨리면 내려가나」}다.
🔴 그리고 \textbf{스팬은 격자가 정한 수}라, 이 순위는 격자 밖으로 일반화되지 않는다 ---
그 한계를 결과와 같은 자리에 적는다.

같은 라벨 위에서 앞 노트의 \textbf{언급 정밀도 추정량도 고친다}.
표본이 등할당 층화인데 헤드라인이 층을 무시한 비율이었고,
제목이 군집인데 구간이 독립 시행을 가정했다.
새 라벨은 \textbf{하나도 붙이지 않았다} --- 바뀐 것은 추정량뿐이다.
\end{abstract}
]

\section{빈 칸은 자가 없어서가 아니라 요인이 없어서 비었다}

여섯 능력 축 가운데 \textbf{「무엇이 capability 를 정하나」}에 붙은 자는
처음부터 분명했다: \textbf{요인 하나만 흔든 짝 실험의 $\Delta \pm SE$}.
그런데도 그 칸은 세 사이클 연속 비었다. 🔴 \textbf{까닭은 자가 아니라 목록이었다} ---
무엇을 흔들 것인지 적어 둔 적이 없으니 흔들 수가 없었다.

같은 공백이 다른 축도 막고 있었다. \textbf{「상태에서 예측으로」} 축은
개선의 \textbf{크기를 못 가른다}는 곳에서 멈춰 있었는데,
크기를 가르려면 \textbf{견줄 다른 크기}가 있어야 한다.
요인 표가 없으면 그 비교 자체가 정의되지 않는다.

그래서 이 노트의 첫 일은 측정이 아니라 \textbf{등록}이다.
학습 파이프라인의 다섯 칸에 각각 요인 하나를 붙이고,
각 요인의 \textbf{격자값을 측정 전에} 사전등록에 박았다.
격자를 뒤에 고치면 실패로 적기로 같이 박았다.

\section{설계 --- 유보를 한 번도 안 건드린다}

밑판은 앞 노트가 쓰던 학습 팔 그대로다. 각 수준은 밑판에서 \textbf{정확히 한 칸만} 다르고,
그것을 산문으로 주장하지 않고 \textbf{기계로 확인해} 산출물에 적는다.

가장 중요한 설계 결정은 \textbf{유보를 손대지 않는 것}이다.
시간으로 자른 유보는 어떤 요인을 흔들어도 바뀌지 않으므로,
모든 수준의 유보 행 지문이 \textbf{하나}여야 한다. 그것도 산출물에 적는다.
$SE$ 는 \textbf{같은 되뽑기}를 두 수준에 동시에 먹이는 짝지은 군집 붓스트랩에서 나온다 ---
수준마다 따로 뽑으면 짝이 깨지고 폭이 부풀기 때문이다.

\section{C4 — 요인 하나만 흔든 짝 실험}

\claimx{요인 다섯을 하나씩만 흔들어 유보 예측 성능의 $\Delta \pm SE$ 를 냈다. 이 격자 안에서의 순위는 특징 집합(architecture) $>$ 학습 창(curriculum) $>$ 원천 혼합(data mixture) $>$ 도메인 가중(objective) $>$ 능형 $\lambda$(optimizer) 다}{한 수준이 밑판과 두 칸 이상 다르거나, 유보 집합이 수준마다 달라지거나, $\Delta$ 에 짝지은 $SE$ 가 안 붙으면 이 주장은 떨어진다. 격자를 넓히면 스팬이 커지므로 순위는 격자 밖으로 일반화되지 않는다.}

밑판은 HPLT 전량 학습($n=37,222$ 행) · $\lambda=1.0$ · 특징 전량6 · 창 전량 · 균등 가중이고, 유보는 모든 수준에서 같다($n=782$ · 게이트 도메인 합 $728$). 밑판 묶음 유보 스피어만은 $0.37815$ 다.

\begin{table}[h]\centering
\begin{tabular}{llll}
\hline
순위 & 요인 & 스팬 & $\Delta \pm SE$ \\
\hline
1 & R2 특징 집합 (architecture) & 0.098849 & -0.098849 ± 0.041639 \\
2 & R3 학습 창 (curriculum) & 0.071424 & -0.071424 ± 0.021517 \\
3 & R5 원천 혼합 (data mixture) & 0.046575 & -0.046575 ± 0.023094 \\
4 & R4 도메인 가중 (objective) & 0.039736 & -0.039736 ± 0.018558 \\
5 & R1 능형 $\lambda$ (optimizer) & 0.017251 & +0.015264 ± 0.012209 \\
\hline
\end{tabular}
\caption{요인 순위. 스팬은 사전등록된 격자 안에서의 최대--최소 차다.}
\end{table}

문턱(사전등록: $|\Delta|$ 가 두 배 $SE$ 이상이고 카드 문턱 이상)을 둘 다 넘은 수준은 6 / 18 이고, 여섯 전부 $\Delta$ 가 음수다. 즉 이 격자에서 잰 것은 성능을 올리는 요인이 아니라 \emph{망가뜨리면 내려가는} 요인이다.

\section{언급 정밀도의 추정량 정정}

\claimx{등할당 층화 표본에서 층을 무시한 비율은 편향 추정량이고, 제목이 군집이면 독립 시행 가정의 구간은 좁다}{층가중 값이 층 무시 값과 같거나, 제목 군집 붓스트랩 구간이 Clopper--Pearson 구간보다 넓지 않으면 이 주장은 떨어진다.}

층가중 개체 정밀도는 $0.771566$ 이고 층을 무시한 값은 $0.75431$ 다. 라벨한 $696$ 행이 제목 $164$ 개에서 나오므로 독립 시행 가정의 Clopper--Pearson 구간은 좁다: 제목 군집 붓스트랩의 95\% 구간은 [0.559859, 0.872423] 이고 설계효과는 $22.5251$ 다. 유효 삼중쌍은 $27499.4$ [19953.9, 31094] 다.


\section{이 순위가 말하지 않는 것}

🔴 \textbf{스팬은 요인의 본질적 크기가 아니다.} 격자를 넓히면 스팬이 커진다.
그래서 이 표는 \textbf{「이 격자 안에서」}의 순위이고, 다음에 필요한 것은
\textbf{격자와 무관한 정규화}(예: 같은 계산 예산에서의 최대 $\Delta$)다.

🔴 \textbf{그리고 문턱을 넘은 수준이 전부 음수라는 사실은 그 자체가 결과다.}
이 격자에서 밑판은 이미 각 요인의 좋은 쪽에 있었고,
흔든 것은 대부분 \textbf{망가뜨리는 방향}이었다. 성능을 올리는 요인을 찾으려면
격자를 밑판의 \textbf{바깥쪽}으로 넓혀야 한다 --- 그것이 다음 사이클의 일이다.

\end{document}
